\documentclass[12pt,a4paper]{article}
\usepackage[utf8]{inputenc}
\usepackage[spanish]{babel}
\usepackage{amsmath}
\usepackage{amsfonts}
\usepackage{amssymb}
\usepackage{graphicx}
\usepackage{geometry}
\usepackage{color}
\usepackage{hyperref}
\usepackage{booktabs}

\geometry{margin=2.5cm}

\title{\textbf{Reporte de Avance: Evaluación, Validación e Integración Final (Iteraciones 3 y 4)}}
\author{Ingeniería en Sistemas - Trabajo Terminal II (TT\_Ratones\_2026)}
\date{\today}

\begin{document}

\maketitle

\section{Introducción}
El presente reporte detalla las actividades y resultados obtenidos durante las Iteraciones 3 y 4 del proyecto. Estas fases se centraron en la validación exhaustiva del sistema, la optimización de hardware para arquitecturas modernas y la consolidación del dashboard interactivo para la entrega final del Trabajo Terminal II.

\section{Iteración 3: Evaluación y Validación}
Esta fase tuvo como objetivo asegurar la robustez científica y técnica del sistema antes de la fase de despliegue final.

\subsection{Mitigación de Riesgos de Hardware}
Un desafío crítico fue la compatibilidad con GPUs de arquitectura Blackwell (RTX 50xx). Se implementaron:
\begin{itemize}
    \item \textbf{Modo de Alta Estabilidad:} Opción para forzar el procesamiento en CPU cuando sea necesario.
    \item \textbf{Gestión de Memoria:} Optimización de fragmentación de memoria en PyTorch (\texttt{PYTORCH\_CUDA\_ALLOC\_CONF}).
\end{itemize}

\subsection{Validación de Componentes}
Se realizaron pruebas funcionales en los siguientes módulos:
\begin{itemize}
    \item \textbf{Ingesta de Video:} Verificación de codecs y recortes temporales precisos.
    \item \textbf{Zonas (ROI):} Validación del factor de escala para mantener la precisión milimétrica independientemente de la resolución de pantalla.
    \item \textbf{Persistencia:} Comprobación de la integridad de datos en PostgreSQL mediante contenedores Docker.
\end{itemize}

\section{Iteración 4: Integración y Dashboard Final}
La Iteración 4 se enfocó en la experiencia del usuario y la generación de resultados automatizados.

\subsection{Dashboard Interactivo (Streamlit)}
Se consolidó una interfaz premium que organiza el flujo de trabajo en 5 módulos lógicos, permitiendo una navegación intuitiva y persistente.
\begin{itemize}
    \item Gestión de sesiones para evitar la pérdida de configuraciones entre páginas.
    \item Sistema de diagnóstico visual de hardware (Docker, DB, GPU) en tiempo real.
\end{itemize}

\subsection{Automatización de SimBA}
Se logró la integración de SimBA mediante scripts de orquestación, permitiendo que el análisis de comportamientos complejos (como el \textit{Grooming}) se realice sin intervención manual en la GUI de SimBA, facilitando el procesamiento por lotes.

\section{Estado Actual y Próximos Pasos}
A la fecha, el sistema es capaz de:
\begin{enumerate}
    \item Detectar y rastrear un roedor con una precisión superior al 99.4\%.
    \item Clasificar comportamientos de ansiedad (Tigmotaxis y EPM).
    \item Persistir toda la información de forma segura en una base de datos relacional.
\end{enumerate}

\textbf{Próximo hito:} Refinamiento del motor de generación de reportes PDF finales para la documentación de resultados experimentales en el TT2.

\section{Conclusión}
Las Iteraciones 3 y 4 han sido completadas satisfactoriamente en sus componentes técnicos fundamentales. El sistema es estable, escalable y cumple con los requisitos técnicos exigidos para el Trabajo Terminal II.

\end{document}
